% Standalone section fragment. Requires amsmath, amssymb, amsthm, hyperref.
\section{A complete finite certificate for the source rotation}
\label{sec:boundary-rotation-certificate}

The construction in this section changes the binary register-program tree,
not the arithmetic algorithm or the searched theory.  Starting with the
6,455-site source of \texttt{candidate-rotation240}, it removes one distinct
internal tree node while retaining a maximum zero run of six.  Its exact
source metrics are
\[
 \begin{array}{c|cc}
 &\text{baseline}&\text{rotated source}\\\hline
 \text{distinct internal nodes}&202&201\\
 \text{distinct leaf instructions}&23&23\\
 \text{total distinct DAG nodes}&225&224\\
 \text{expanded leaf occurrences}&6455&6455\\
 \text{maximum leaf-path length}&51&48\\
 \text{maximum consecutive zero bits}&6&6.
 \end{array}
\]
The reduction is certified for every natural register vector.  A binary
Turing-machine simulation theorem is a separate obligation: context-sensitive
decoder folds need not survive a correct source reassociation.

\subsection{Exact source semantics, including path boundaries}

Let $T$ be a finite full binary tree.  Its leaves carry the instructions
$\mathrm{CI}(r)$, $\mathrm{CD}(r)$, $\mathrm{CB}(k)$,
$\mathrm{CC}(k)$, or $\mathrm{CH}$.  Number the leaf occurrences
$0,\ldots,m-1$ in left-to-right order and write $p_i$ for leaf $i$'s path.
For any finite binary word $P$, define $\operatorname{loc}_T(P)$ by
following $P0^\omega$ from the root, stopping at the first leaf.  Any bits
remaining after the encountered leaf are ignored.  In particular,
$\operatorname{loc}_T(\epsilon)=0$.

Write $\operatorname{inc}(p)$ for fixed-length binary increment, with the
convention
\[
 \operatorname{inc}(1^d)=0^d,
 \qquad \operatorname{inc}(\epsilon)=\epsilon.
\]
Write $\operatorname{strip}_b(k,p)$ for removal of a suffix through the
$k$th occurrence of bit $b$, counting from the right.  If the word is
exhausted first, the result is $\epsilon$; for $k=0$, it is $p$.  Set
\[
 \operatorname{brk}(k,p)=
 \operatorname{inc}(\operatorname{strip}_0(k,p)),
 \qquad
 \operatorname{cont}(k,p)=\operatorname{strip}_1(k,p).
\]
Every instruction operates on the \emph{encountered leaf path} $p_i$,
not the untruncated input word $P$.  This is exactly the role of
\texttt{firstn d (P ++ repeat false d)} in the upstream definition of
\texttt{rm\_step}.  These conventions reproduce the explicit definitions
in \href{https://github.com/jeff87654/busy-beaver/blob/zf-254/ZF/m254/coq/tm/RM.v}
{\texttt{RM.v}} and its
\href{https://github.com/jeff87654/busy-beaver/blob/zf-254/ZF/m254/tools/prog/rmpath.py}
{Python path-semantics mirror}.

The complete finite control descriptor $G_T(i)$ is
\[
\begin{array}{ll}
 (\mathrm{CI},r,\operatorname{loc}_T(\operatorname{inc}(p_i)))
       &\text{for }\mathrm{CI}(r),\\
 (\mathrm{CD},r,\operatorname{loc}_T(\operatorname{inc}(p_i)),
             \operatorname{loc}_T(\operatorname{brk}(1,p_i)))
       &\text{for }\mathrm{CD}(r),\\
 (\mathrm{J},\operatorname{loc}_T(\operatorname{brk}(k,p_i)))
       &\text{for }\mathrm{CB}(k),\\
 (\mathrm{J},\operatorname{loc}_T(\operatorname{cont}(k,p_i)))
       &\text{for }\mathrm{CC}(k),\\
 (\mathrm{H})&\text{for }\mathrm{CH}.
\end{array}
\]
Both branches of every decrement occur in this descriptor.  No assertion
that a register is usually zero, or that a branch is not reached in a
finite experiment, enters its definition.  Representing break and continue
by the common symbol $\mathrm J$ permits exact rebasing of explicit control
instructions, when needed; neither modifies a register.

\paragraph{Finite-control isomorphism lemma.}
Suppose $T,U$ are finite full binary source trees with the same number of
leaves and $G_T(i)=G_U(i)$ for every leaf index $i$.  Then their path
semantics are strongly bisimilar from the initial empty path, for every
initial register vector.  Their register vectors agree at every matched
instruction, and they reach their halt instructions at the same source
step exactly when either reaches one.

\begin{proof}
Relate configurations $(P,R)$ and $(Q,S)$ precisely when $R=S$ and
$\operatorname{loc}_T(P)=\operatorname{loc}_U(Q)$.  Initial empty paths
are related because both locate leaf zero.  Assume related configurations
locate index $i$.

For an increment descriptor, the same register is incremented on both
sides, and the descriptor supplies the same successor index.  For a
decrement descriptor, equality of the register vectors makes the zero
test agree.  In the zero case both vectors are unchanged and the two zero
successor indices agree.  In the positive case both decrement the same
positive register and the two nonzero successors agree.  For a jump
descriptor both vectors remain unchanged and the successor indices
agree.  For a halt descriptor both sides halt.  Each case therefore
preserves the relation or halts on both sides.

This argument uses the exact zero-padded walk and the actually consumed
leaf path.  It covers root rollover, an empty path after stripping, paths
shorter than the next leaf, and paths longer than the leaf encountered.
Induction on the number of source steps proves the claim.  Full finite
trees have no missing branch or cyclic internal edge, so a walk cannot
become stuck; consequently halting here means reaching $\mathrm{CH}$,
not exhausting a traversal fuel bound.
\end{proof}

The lemma is deliberately independent of an associativity heuristic.
An arbitrary rotation of a source tree can change a decrement's implicit
break target and be unsound.  A rotation is accepted here only when its
entire resulting descriptor array agrees with the pinned baseline array.

\subsection{Thirteen replayable stages}

A left rotation means $(A,(B,C))\mapsto((A,B),C)$; a right rotation is its
inverse.  Paths in the certificate always refer to the tree immediately
before that stage.  Multiple paths in one stage are prefix-incomparable,
so they specify simultaneous disjoint rewrites.

Stage 1 applies the same left rotation at all 13 occurrences of the source
fragment identified in the search as node 69.  Its three closed chunks are
\texttt{PUSH}, \texttt{copy-param1}, and
\texttt{PUSH;copy-cache}.  It reuses the existing paired chunk
\texttt{PUSH;copy-param1} and removes one distinct internal source subtree.
Stages 2--11 are the closed-chunk reassociations recovered from the saved
source search.  The remaining stages cross the final explicit break in
the AX command; they are exactly the two moves absent from a rule that
requires all three chunks of a rotation to be closed.

\begin{center}
\begin{tabular}{rcl}
Stage&Direction&Path (before that stage)\\\hline
1&left&13 explicit paths in \texttt{rotation\_certificate.json}\\
2&left&\texttt{1110100100101011010101000}\\
3&left&\texttt{111010010010101101010}\\
4&right&\texttt{11101001001010110101}\\
5&left&\texttt{11101001001010110101010010}\\
6&right&\texttt{1110100100101011010101001}\\
7&right&\texttt{11101001001010110101010011}\\
8&left&\texttt{111010010010101101010100111}\\
9&left&\texttt{11101001001010110101010000}\\
10&right&\texttt{1110100100101011010101000}\\
11&left&\texttt{111010010010101101010100111}\\
12&right&\texttt{11101001001}\\
13&right&\texttt{111010010011}
\end{tabular}
\end{center}

There are 13 stages and 25 individual tree rotations, with no change to
any leaf instruction or explicit break/continue operand.  Stage 12 turns
$((\mathrm{GEN},\mathrm{AXrest}),\mathrm{CB}(1))$ into
$(\mathrm{GEN},(\mathrm{AXrest},\mathrm{CB}(1)))$.  It removes the sole
remaining excessive zero run at leaf 40, but transfers it to leaf 121,
the next parameter-unpair decrement.  Stage 13 reassociates the two
closed portions of that remainder across the same trailing break and
removes the transferred excessive run as well.  The first eleven stages
end with maximum zero run seven and one offending leaf; the last stage
ends with maximum zero run six and none.  Both final moves leave the
number of distinct DAG nodes at 224.

\subsection{Portable independent verification}

The following five files suffice to replay and verify the entire source
construction, using only the Python standard library:
\begin{itemize}
\item \texttt{baseline\_sites.json}: the exact pinned source input;
\item \texttt{rotation\_certificate.json}: all 13 stages and all paths;
\item \texttt{boundary224\_sites.json}: the claimed final source;
\item \texttt{check\_rotation\_certificate.py}: the independent checker;
\item \texttt{rotation\_certificate\_results.json}: the recorded result.
\end{itemize}
The compact \texttt{boundary224\_dag.json} additionally provides the 224
distinct source nodes.  Its metric counts include the 23 distinct leaf
instructions and must not be mistaken for a binary Turing-machine state
count.

Run
\begin{quote}\ttfamily
python3 check\_rotation\_certificate.py
\end{quote}
in the artifact directory.  The checker verifies the baseline byte hash,
validates a finite full binary source tree, replays every structural
rotation, checks the ordered instruction word, and compares the complete
control descriptor arrays after each stage.  Each stage checks 8,385
successor/guard edges and the unique halt descriptor.  It then compares
the replayed source against the claimed final source and recomputes all
DAG and path metrics.  It has no register-value cutoff, execution fuel,
heuristic reachability test, repository import, or solver dependency.
All checks use explicit exceptions rather than Python assertions; running
with \texttt{python3 -O} does not disable any certificate check.

For every generated successor word, the checker also verifies that any
unused suffix after the encountered leaf contains only zeros.  This is a
stronger canonicality check than the control-isomorphism lemma itself
requires.  Together with the initial empty word, it proves by induction the precise
canonicality identity
\[ P0^\omega=p_{\operatorname{loc}_T(P)}0^\omega. \]
This covers a word shorter than its encountered leaf path as well as a word
with redundant trailing zeros.  Thus no nonzero discarded tail can silently
alter the path convention used by a later machine interface.

The canonical normalized control-graph SHA-256 is
\begin{quote}\ttfamily\small
ae7f12f37306b86e5f70469210fd2a8d\\
68b6f9df094bf659abc6b4a56714f1fd.
\end{quote}
The exact baseline source SHA-256 is
\begin{quote}\ttfamily\small
74b1e7f73cc3ffe7761e36ab2d57a26e\\
0ba34d3ac624f2472799e6ac41d5a51d.
\end{quote}
The exact final source file SHA-256 is
\begin{quote}\ttfamily\small
7fd152d1f5371d2a561d4ad1870d2416\\
83b010ce7b7a27399b95363a4537cfb6.
\end{quote}
The line breaks in these displays are purely typographic.  The result
JSON records the hashes as uninterrupted strings.  As an additional
cross-check, the upstream segment-bisimulation
generator independently produced a 2,987-entry certificate between the
baseline and the final source.  The portable checker above does not
depend on that generator or certificate.

These are executable finite certificates plus an ordinary mathematical
proof of their significance.  They are not, by themselves, a Lean/Rocq
kernel proof or a serialized derivation in arithmetic.  They establish source equivalence.  Sections~\ref{sec:boundary240} and
\ref{sec:fold239} supply the separate instruction-path decoding, context
repair, blank initialization, and unbounded binary simulation obligations
for the final 239-state construction.
